\section{Institutional Background}
\label{sec:background}

\subsection{From Incurred Losses to Expected Credit Losses}
\label{sec:background:transition}

Under the historical incurred loss model (ILM), U.S.\ banks recognized
credit losses only when a loss was probable and could be reasonably
estimated, producing loan-loss provisions that lagged realized defaults
\citep{beatty2011delays} and reduced the informativeness of bank
disclosures \citep{bushman2012accounting, bushman2015delayed,
bischof2021accounting}.  Following
the 2007--2009 crisis, FASB issued ASU 2016-13, introducing the Current
Expected Credit Loss (CECL) model, codified in ASC~326
\citep{FASB2016}: banks estimate and reserve for \emph{lifetime}
expected credit losses at origination and at each reporting date.  The
required allowance is a direct function of loan-portfolio credit-risk
composition: longer maturities, weaker borrowers, and thinner collateral
require larger allowances than the ILM reserve.  Forcing credit risk
into a reported number in this way is what makes the standard
informative to outsiders; the same logic underlies the finding that
fair-value reporting of financial instruments strengthens the observed
association between bank leverage and credit risk
\citep{blankespoor2013fair}.  Recent evidence finds
CECL adoption made provisions timelier and more reflective of future
local economic conditions \citep{kim2026current}, raised rates on
longer-maturity consumer credit \citep{granja2025current}, tightened
lending in the COVID-19 recession \citep{chen2022does}, and produced
Day-One adjustments that predict future net charge-offs and are
reflected in equity prices \citep{gee2022decision}.

\subsection{The Day-One Adjustment and Treatment Variable}
\label{sec:background:dayone}

When a bank transitions from ILM to CECL, it must reclassify its entire
existing loan portfolio under the new standard on a single date.  The
difference between the CECL-required allowance and the legacy ILM allowance
(the Day-One adjustment) is recorded as a cumulative-effect charge to
retained earnings on the opening balance sheet of the adoption quarter.  For
bank $i$ with adoption quarter $t^{*}$:
\begin{equation}
  \Delta_i
  \;\equiv\;
  \mathrm{ALLL}^{\mathrm{CECL}}_{i,\,t^{*}}
  -
  \mathrm{ALLL}^{\mathrm{ILM}}_{i,\,t^{*}-1},
  \label{eq:dayone}
\end{equation}
where $\mathrm{ALLL}$ is the allowance for loan and lease losses on the Call
Report.  The Day-One adjustment varies widely: a large share of community
banks recorded a negligible adjustment, a subset recorded positive values
where CECL required a materially larger allowance, and a smaller number
recorded negative values.

Where positive, the adjustment reduces reported book equity and, absent
regulatory accommodation, risk-weighted capital ratios by an equivalent
amount.  It is a pure accounting event: cash flows and borrower repayment
terms are unchanged.  What the adjustment reveals, in a single mandatory
disclosure, is the gap between book equity as previously reported and book
equity adjusted for the bank's forward-looking credit-loss exposure.

The primary treatment scales the Day-One charge by beginning-of-period
equity:
\begin{equation}
  \texttt{CECL Adj}_{i}
  \;=\;
  \frac{\Delta_i}{\texttt{Equity}_{i,\,t-1^{*}}}
  \times 100,
  \label{eq:treatment}
\end{equation}
expressed in percentage points.  For binary specifications I construct
\begin{equation}
  \texttt{High CECL}_{i}
  \;=\;
  \mathbf{1}\!\left\{
    \texttt{CECL Adj}_{i} \geq \tau
  \right\},
  \label{eq:binary}
\end{equation}
with $\tau = 2.5$ percentage points (approximately the 95th percentile).
Equity scaling is the baseline because it measures the fraction of the
loss-absorbing buffer that the new standard revealed to be already
committed against expected credit losses.

\subsection{Is the Adjustment Predetermined?}
\label{sec:background:predetermination}

The 2023Q1 adoption date was set by regulation and cannot have been
chosen to coincide with favorable deposit conditions.  The size of the
Day-One adjustment reflects the cumulative composition of the loan
portfolio, fixed by prior underwriting: long maturities, weak borrower
credit, thin collateral.  The standard was public from June 2016, but
the \emph{magnitude} of any given bank's Day-One charge was not
knowable to outsiders in advance; anticipation of the event date
without knowledge of intensity does not bias a
difference-in-differences design, and outside anticipation of intensity
would appear as pre-adoption divergence in depositor-facing outcomes,
which I test for directly.

Two components of the Day-One magnitude deserve separate treatment.
The first is legacy loan composition, which was fixed by prior
underwriting and is uncorrelated with 2023 deposit dynamics under any
reasonable assumption.  The second is the forecast and modeling
component required by ASC~326: banks combine historical loss experience
with reasonable-and-supportable forward-looking forecasts of
macroeconomic conditions and choose a lifetime-loss estimation
approach.  This component is a contemporaneous choice at adoption and
is in principle correlated with 2022--2023 macro conditions.  This
choice does not threaten the design.  Uninsured and insured time
deposits at the same bank face the same modeling choice, so it biases
the composition comparison only if it moves the two deposit types
apart through some channel other than the disclosure itself.

\subsection{What the Disclosure Does, and Does Not, Change}
\label{sec:background:mechanism}

The premise of an information shock is that bank credit risk is hard
to observe from outside.  Banks are the canonical opaque industry:
rating agencies split over bank bonds far more often than over other
issuers, and the disagreement traces precisely to loans
\citep{morgan2002rating}.  Mandated supervisory disclosure
demonstrably closes part of that gap, but the demonstrations are for
capital-market audiences: stress-test results move stock prices and
trading volume, most strongly at levered and risky banks
\citep{flannery2017stress}.  Whether a mandated credit-risk disclosure
reaches depositors is the question this paper takes up.

Prior depositor-discipline tests typically exploit realized distress in
which new information arrives simultaneously with deterioration in cash
flows or liquidity \citep{martinez2001depositors, iyer2016tale}.  The CECL Day-One adjustment separates the two: it
records a pure accounting charge that does not trigger a cash transfer,
alter borrower repayment terms, reduce bank liquidity, or change the
risk of the underlying loans.  A depositor who withdraws or demands a
higher rate responds to newly disclosed credit-risk information, not to
concurrent worsening of the bank's actual condition.

\subsection{Adoption Timeline and Market Validation}
\label{sec:background:timeline}

Large SEC-filing registrants adopted CECL in 2020Q1; non-public
entities, the community banks in the main sample, adopted in 2023Q1.
Call Reports are posted publicly by the FFIEC and FDIC roughly thirty
days after quarter-end, so 2023Q1 Day-One adjustments became publicly
observable in late April and May 2023, after the March 2023 failures of
Silicon Valley Bank and Signature Bank.  Depositors reacting to the
disclosure did so against an information backdrop in which the March run
had already occurred.

The 2020 large-bank cohort shows that market participants treat the
Day-One number as material news.  Using FR~Y-9C financial statements,
CRSP equity market capitalization (linked via the CRSP--FRB crosswalk),
and Bloomberg five-year CDS spreads,
Figure~\ref{fig:es_large_banks_market_cds} reports monthly event
studies around the 2020Q1 adoption for the SEC-filing cohort,
interacting event-time dummies with each holding company's CECL
adjustment scaled by equity.  Equity markets respond immediately:
relative to the month before adoption, the coefficient on normalized
market capitalization drops to $-0.947$ per unit of CECL-to-equity at
the adoption month ($p = 0.006$) and reaches a trough of $-2.771$
($p < 0.001$) two months later; a ten-percentage-point higher
adjustment is associated with roughly 28~percent lower market
capitalization.  Credit markets price the same information with a lag:
five-year CDS spreads for the 103 banks with active contracts widen by
$190.8$~basis points per unit of CECL-to-equity four months after
adoption ($p = 0.047$), about 19~basis points per ten percentage
points of adjustment.  Pre-adoption coefficients in both markets are
flat.

I use the 2020 cohort only for this signal-validation exercise, not
for the deposit tests, for two reasons.  The cohort is small: roughly
one hundred publicly traded holding companies, against more than 4,000
community banks.  And its adoption quarter coincides with the COVID-19
onset, when emergency liquidity programs and precautionary corporate
drawdowns flooded banks with deposits; any disciplinary outflow would
be swamped by that inflow, and no clean counterfactual exists within
the cohort.  Equity and CDS prices, by contrast, can absorb the
disclosure even in a turbulent quarter because the event studies
compare across banks within the same month.
